% Draft for IEEE ICORT 2027 (IEEEtran conference). Numbers marked \tbd{} come from
% scripts/run_efficiency.sh and are filled in once that run finishes.
\documentclass[conference]{IEEEtran}
\usepackage{amsmath,amssymb,booktabs,graphicx,xcolor,multirow}
\usepackage[hidelinks]{hyperref}
\newcommand{\tbd}[1]{\textcolor{red}{[#1]}}

\begin{document}

\title{Train as Experts, Deploy as Dense: Sequence-Level Mixture-of-LoRA Adaptation of a 4-bit Vision-Language Model for Edge Devices}

\author{\IEEEauthorblockN{\tbd{Author names}}
\IEEEauthorblockA{\tbd{Department, Institute}\\ \tbd{email}}}

\maketitle

\begin{abstract}
Vision-language models (VLMs) are attractive for on-device document, chart and scene understanding, but
their memory footprint and the cost of adapting them to several tasks limit use on resource-constrained
hardware. We adapt Qwen2-VL-2B, quantised to 4-bit NF4, to three tasks (document VQA, chart QA and
visual spatial reasoning) with a mixture of LoRA experts (MoE-LoRA) trained on a single 16\,GB T4 GPU
with knowledge distillation from the full-precision model. We compare token-level routing with a
sequence-level router that makes one expert choice per request from the prompt. All adapters improve
the 4-bit model by 6--7 points on held-out test data. Sequence-level MoE-LoRA matches a dense LoRA of
the same rank within noise ($-0.7$ points, 95\% CI $[-1.7, +0.4]$, three seeds) while its experts become
strongly task-aligned without task labels: the mutual information between task and expert is about
$10\times$ that of token-level routing in every seed. Because each request uses one expert, the
experts fold into standard per-task checkpoints with no accuracy loss. On a 4\,GB GTX~1650 laptop GPU
the folded model runs at exactly the base model's cost (2.05\,GB peak memory, 18.8 tokens/s), whereas
keeping the router costs 11\% more memory and 24\% slower decoding. \tbd{Efficiency sentence: accuracy
of 4-bit + adapters at 256\,px vs the fp16 model at 512\,px, memory and image-token savings.}
\end{abstract}

\begin{IEEEkeywords}
vision-language models, mixture of experts, LoRA, quantization, edge deployment, knowledge distillation
\end{IEEEkeywords}

\section{Introduction}
Compact VLMs such as Qwen2-VL-2B~\cite{wang2024qwen2vl} can read documents and charts and reason about
scenes, which makes them candidates for offline use on laptops, embedded boards and field equipment.
Two obstacles remain. First, the full-precision model needs about 4.4\,GB for its weights alone, more
than many edge GPUs offer. Second, one device often serves several distinct tasks, and adapting one
model to all of them either costs accuracy or multiplies the number of models to store.

Parameter-efficient fine-tuning with low-rank adapters (LoRA)~\cite{hu2022lora} on a 4-bit quantised
base (QLoRA)~\cite{dettmers2023qlora} addresses the first obstacle during training. Mixtures of LoRA
experts~\cite{zadouri2023molora,dou2023loramoe} address the second by giving the model several
adapters and a router, in the spirit of sparse mixture-of-experts layers~\cite{shazeer2017moe,fedus2022switch}.
However, routed experts add memory and latency at inference, which is exactly what an edge device cannot
afford, and token-level routers rarely learn experts that correspond to recognisable tasks.

We study whether MoE-LoRA can be trained for multi-task capacity and still be deployed at the cost of
the base model. Our contributions are:
\begin{itemize}
\item A shared-base MoE-LoRA layer for the MLP blocks of a 4-bit VLM, with a \emph{sequence-level}
router that selects one expert per request from the pooled prompt text, trained with sparse
top-$k$ knowledge distillation from the full-precision model on a single T4 GPU.
\item A controlled three-seed comparison against dense LoRA (rank-matched and parameter-matched) and
token-level routing on held-out test data, with paired bootstrap tests.
\item An analysis showing that sequence-level routing yields task-aligned experts without labels
($\approx$0.48 vs 0.05 bits of task--expert mutual information), and that these experts fold
losslessly into standard per-task checkpoints.
\item Measurements on a 4\,GB consumer laptop GPU showing that folded checkpoints have zero overhead
over the base model, including a practical finding for GPUs without tensor cores.
\end{itemize}

\section{Related Work}
\textbf{Parameter-efficient adaptation.} LoRA~\cite{hu2022lora} learns low-rank updates to frozen
weights; QLoRA~\cite{dettmers2023qlora} trains them on top of a 4-bit NF4 base, which we adopt.
\textbf{Mixture of experts.} Sparsely gated MoE layers~\cite{shazeer2017moe} and Switch
Transformers~\cite{fedus2022switch} route tokens to experts with load-balancing losses; router
$z$-losses~\cite{zoph2022stmoe} stabilise training. Mixtures of LoRA experts apply the idea to
adapters~\cite{zadouri2023molora,dou2023loramoe}. These works route at token level and keep the router
at inference. We route at sequence level and fold experts away for deployment.
\textbf{Distillation.} We use temperature-scaled knowledge distillation~\cite{hinton2015distilling} from
the unquantised model to the 4-bit student.

\section{Method}
\subsection{Shared-base MoE-LoRA layer}
Each MLP block of the language model, $y=W_d\,(\sigma(W_g x)\odot W_u x)$, keeps its frozen 4-bit
weights and receives $E$ LoRA experts on each of the three projections:
\begin{equation}
W_p' x = W_p x + \sum_{e=1}^{E} g_e(x)\, s\, B_{p,e} A_{p,e}\, x,\quad p\in\{g,u,d\},
\end{equation}
with rank $r$, scale $s=\alpha/r$ and gate vector $g(x)$. The base projection is computed once and never
scaled, so the layer is exactly the base model when all $B=0$. The $E$ adapters of a projection are
stored stacked, and the gate is applied as a block mask on the rank dimension, so the cost equals one
LoRA of rank $Er$ without a loop over experts. For the selected top-$k$ experts the gate is
$g_e=(E/k)\,p_e$, where $p$ is the router's softmax output; a dense LoRA baseline is the special case $E=1$.

\subsection{Token-level vs sequence-level routing}
A token-level router computes $p$ from each token's hidden state, as in standard MoE. Our
sequence-level router computes one decision per request: it mean-pools the layer input over the
prompt's \emph{text} tokens (image tokens excluded), applies the router once, and uses the chosen
expert for every token of the sequence. During generation the decision made in the prefill pass is
reused for the decoded tokens, so the answer cannot influence the choice. Because a single decision per
sample makes per-batch load statistics meaningless at batch size 1, the load-balancing loss
$E\sum_e \bar u_e \bar p_e$ uses an exponential moving average $\bar u$ of expert usage (decay 0.99);
a router $z$-loss~\cite{zoph2022stmoe} is added.

\subsection{Training objective}
The loss is $\mathcal{L}=\mathcal{L}_{CE}+\lambda_{KD}\mathcal{L}_{KD}+\lambda_{lb}\mathcal{L}_{lb}+\lambda_z\mathcal{L}_z$,
computed on answer tokens only. The teacher is the same model in fp16 without quantisation. We cache its
top-50 probabilities at temperature $T=2$ for every answer token. Because at $T=2$ the top-50 tokens hold
on average only 44\% of the probability mass, we do not renormalise them; instead the remaining mass forms
one ``tail'' bucket, and $\mathcal{L}_{KD}$ is the exact KL divergence between the $(k{+}1)$-bucket
teacher and student distributions, scaled by $T^2$. It is zero for an identical student.

\subsection{Folding experts for deployment}
In \emph{fixed} mode the gate is a constant vector $\bar g$ and each projection's update is the dense
matrix $\sum_e \bar g_e s B_e A_e$, which is added to the unquantised base weights and saved as a standard
Qwen2-VL checkpoint. A \emph{task profile} uses, for each layer, the expert that the router chose most often
for that task on the development set, with its mean gate value; a \emph{full merge} uses the expected gate
over all development samples. With sequence-level routing a request uses exactly one expert per layer, so a
task profile reproduces the routed computation whenever the router's choice agrees with the profile.

\section{Experimental Setup}
\textbf{Data.} 400/50/200 training/development/test samples per task from DocVQA~\cite{mathew2021docvqa}
(document\_ocr), ChartQA~\cite{masry2022chartqa} (chart\_qa) and VSR~\cite{liu2023vsr} (spatial\_reasoning),
1200/150/600 in total. Metrics: ANLS~\cite{biten2019stvqa} for documents, relaxed accuracy for charts and
accuracy on True/False for VSR; we report their macro average. The development set is used for epoch
selection only; all main results are on the test set.

\textbf{Model and training.} Qwen2-VL-2B-Instruct with the language model in 4-bit NF4 (double
quantisation, fp16 compute). Images are resized by the processor to at most 401{,}408 pixels
(``512\,px''), chosen by a resolution sweep on the base model (dev macro 70.5/75.9/76.5 at
256/512/768\,px). Variants: \emph{dense16} (LoRA rank 16, 14.1M parameters), \emph{dense64} (rank 64,
56.4M, parameter-matched to the MoE), \emph{moe} (4 experts of rank 16, top-1, token-level routing, 56.6M)
and \emph{moe\_seq} (the same with sequence-level routing). AdamW, learning rate $2\times10^{-4}$
($10^{-3}$ for routers), cosine schedule with 5\% warm-up, batch 1 with 8-step accumulation, 2 epochs,
gradient checkpointing; $\lambda_{KD}=1$, $\lambda_{lb}=0.01$, $\lambda_z=10^{-3}$. One run takes 64--69
minutes on one T4. Seeds 0--2 for dense16 and moe\_seq, 0--1 for moe, 0 for dense64; repeated runs with
the same seed reproduce identical scores. Significance: paired bootstrap over test samples (stratified by
task, 10{,}000 resamples, per-sample scores averaged over seeds).

\textbf{Hardware.} Training and accuracy evaluation: NVIDIA T4 (16\,GB). Edge measurements: HP Victus laptop,
GTX~1650 (4\,GB, no tensor cores), Intel i5-12450H, 8\,GB RAM. Greedy decoding, batch 1.

\section{Results}
\subsection{Accuracy}
\begin{table}[t]
\centering
\caption{Held-out test accuracy (\%, mean $\pm$ s.d. over seeds). $\Delta$: difference to dense16 with 95\% paired bootstrap CI.}
\label{tab:main}
\setlength{\tabcolsep}{3pt}
\begin{tabular}{lcccccc}
\toprule
Model & Seeds & Chart & Doc & Spatial & Macro & $\Delta$ [95\% CI] \\
\midrule
Base (4-bit) & -- & 72.0 & 82.2 & 60.5 & 71.6 & $-6.8$ [$-9.7$, $-4.1$] \\
dense16 & 3 & 73.5\scriptsize{$\pm$1.3} & 83.5\scriptsize{$\pm$0.8} & 78.2\scriptsize{$\pm$2.1} & \textbf{78.4}\scriptsize{$\pm$0.7} & -- \\
dense64 & 1 & 76.5 & 83.0 & 75.0 & 78.2 & $-0.2$ [$-2.4$, $+2.0$] \\
moe (token) & 2 & 74.5\scriptsize{$\pm$0.7} & 83.3\scriptsize{$\pm$1.8} & 73.2\scriptsize{$\pm$3.9} & 77.0\scriptsize{$\pm$0.9} & $-1.4$ [$-2.7$, $0.0$] \\
moe\_seq & 3 & 73.0\scriptsize{$\pm$1.7} & 84.4\scriptsize{$\pm$0.5} & 75.8\scriptsize{$\pm$0.8} & 77.7\scriptsize{$\pm$0.9} & $-0.7$ [$-1.7$, $+0.4$] \\
\bottomrule
\end{tabular}
\end{table}

Table~\ref{tab:main} shows that every adapter improves the 4-bit base model by 6--7 points. The
sequence-level MoE is statistically indistinguishable from dense LoRA, while token-level routing is
measurably worse ($-1.4$, CI upper bound 0.0). A four times larger dense adapter (dense64) brings no gain,
so adapter capacity does not limit accuracy in this setting. On the development set the ordering of
dense16 and moe\_seq was reversed (84.6 vs 84.8), a difference that did not survive on unseen data; we
therefore make no claim that MoE improves accuracy.

\subsection{Expert specialisation}
\begin{table}[t]
\centering
\caption{Mutual information between task and chosen expert (bits; maximum $\log_2 3=1.585$), development set, mean over the 28 layers.}
\label{tab:mi}
\begin{tabular}{lccc}
\toprule
Routing & Seed 0 & Seed 1 & Seed 2 \\
\midrule
Token-level (moe) & 0.051 & 0.046 & -- \\
Sequence-level (moe\_seq) & 0.526 & 0.469 & 0.436 \\
\bottomrule
\end{tabular}
\end{table}

Table~\ref{tab:mi} measures how much the router's choice reveals about the task. Token-level routing
spreads every task over all experts. Sequence-level routing yields roughly ten times more task information
in every seed, although no task label is ever given to the model. In the most specialised layer of seed 0,
charts, documents and spatial questions go to three different experts in 100\%, 84\% and 94\% of cases.
Charts and documents often share an expert in other layers, consistent with both requiring text reading.
Which layers specialise varies between seeds, and some layers collapse to a single expert (usage entropy 0),
which then acts as a dense LoRA.

\subsection{Folding experts into standard checkpoints}
For moe\_seq, per-task profile checkpoints score 76.9, 78.2 and 77.5 on test for seeds 0--2, against
76.9, 78.7 and 77.6 for the routed models (mean $-0.2$). A single full merge without task information loses
0.7 points on the development set (83.9 vs 84.6). Because each profile is a standard Qwen2-VL checkpoint,
it needs no custom inference code.

\subsection{Edge deployment}
\begin{table}[t]
\centering
\caption{GTX 1650 (4\,GB) laptop GPU, 4-bit weights, fp32 compute, batch 1.}
\label{tab:edge}
\setlength{\tabcolsep}{4pt}
\begin{tabular}{lccc}
\toprule
Model & Peak VRAM & Latency & Decode \\
\midrule
Base & 2054\,MiB & 3.09\,s & 18.8 tok/s \\
moe\_seq, routed & 2278\,MiB & 3.25\,s & 14.2 tok/s \\
moe\_seq, folded & 2054\,MiB & 3.08\,s & 18.8 tok/s \\
dense16, folded & 2054\,MiB & 3.08\,s & 18.7 tok/s \\
\bottomrule
\end{tabular}
\end{table}

Table~\ref{tab:edge} shows that a folded checkpoint has exactly the memory and speed of the base model,
while keeping the router costs 11\% more memory and 24\% slower decoding. About 2.9\,s of the 3.1\,s
latency is the prefill pass, dominated by the vision encoder, so the number of image tokens is the main
lever for latency. The fp16 model cannot be loaded on this GPU: its 4.4\,GB of weights exceed the 4\,GB of VRAM. On the GTX~1650,
which has no tensor cores, fp16 matrix multiplication is about five times slower than fp32; computing in
fp32 with 4-bit weights reduces time to first token from 12.8\,s to 3.0\,s at the same accuracy for
0.5\,GB more memory.

\subsection{Accuracy versus memory and resolution}
\tbd{From run\_efficiency.sh: test accuracy, weights and peak memory of the fp16 model at 512\,px; 4-bit
base at 256/512\,px; moe\_seq and dense16 trained at 256\,px (3 seeds), with paired bootstrap against the
fp16 model. Claim to state only if supported: matching the fp16 model with $\approx$3$\times$ less weight
memory and $\approx$45\% fewer image tokens.}

\section{Conclusion}
Sequence-level MoE-LoRA lets a 4-bit VLM be trained with several task experts on one T4 GPU and deployed
as ordinary per-task checkpoints with no accuracy loss and no runtime overhead on a 4\,GB laptop GPU.
It matches, but does not exceed, a dense LoRA in accuracy; its benefit is modularity and interpretable,
label-free task specialisation at zero deployment cost. Token-level routing, by contrast, learns little
task structure and is less accurate. Limitations: three tasks and small training sets (400 per task),
a single 2B backbone, and task profiles that assume the task is known at deployment.

\bibliographystyle{IEEEtran}
\bibliography{refs}
\end{document}
